\documentclass[11pt]{article}

\usepackage[a4paper,inner=1.2in,outer=1.2in,top=0.9in,bottom=1.0in]{geometry}
\usepackage{newtxtext,newtxmath}
\usepackage{microtype}
\usepackage{amsmath,amssymb}
\usepackage{booktabs}
\usepackage{graphicx}
\usepackage{siunitx}
\usepackage{hyperref}

\pagestyle{empty}

\begin{document}


\centering
{\large \textbf{State of Differences}}\\[0.6cm]
{\large \textbf{\textit{DenoiseRep\textsuperscript{++}}: Enhanced Denoising Model for Representation Learning}}\\[0.6cm]
{\large Submitted to IEEE Transactions on Pattern Analysis and Machine Intelligence}\\[0.6cm]

\raggedright

\hspace*{1em}
This work extends our NeurIPS 2024 version~\cite{xu2024denoiserep} with \textbf{substantial methodological, experimental, and performance enhancements}, forming a significantly strengthened submission suitable for TPAMI.\\[1.2em]


% --------------------------------
% 1. 方法改进
% --------------------------------
\textbf{1. Methodological Improvements}

\begin{itemize}

    \item \textbf{Teacher–Student Distillation for Denoising Module:}  
    The constraint of ensuring a computation-free inference necessitates a highly compact denoising module. However, we identify that such a lightweight structure suffers from limited capacity, leading to underfitting of the complex high-dimensional noise distribution in the NeurIPS version. To resolve this capacity-efficiency trade-off, we introduce a teacher–student distillation mechanism in which a stronger teacher denoiser supervises 
    a compact student. This provides a smoother and more accurate noise estimation signal (see Section D of the Methods, Eqs. (11) and (12)):
    \[
    \mathcal{L}_{distill} = \|\hat{\epsilon}_S(X_t,t) - \hat{\epsilon}_T(X_t,t)\|_2^2.
    \]

    \textit{Evidence from metrics:}  
    Compared with the baseline \textit{DenoiseRep} module, distillation yields more discriminative feature representations, characterized by sharper class boundaries and reduced ambiguity. Specifically, Top-1 accuracy on ImageNet 
    improves from 82.13\% (\textit{DenoiseRep}) to 82.30\% (\textbf{+0.17\% over \textit{DenoiseRep}}), and MSMT17 
    ReID mAP increases from 67.33\% to 67.73\% (\textbf{+0.40\% over \textit{DenoiseRep}}).

    \item \textbf{Multi-Step Denoising Fusion:}  
    Inspired by progressive sampling in diffusion models, we mathematically merge the parameters of multiple denoising steps into a single linear transformation via re-parameterization. This effectively compresses a multi-step denoising trajectory into a single operation. This enhances feature refinement without increasing model size or inference cost. While explicitly executing the denoising module (the non-fused alternative) increases inference latency by approximately 15\% (as detailed in Table V), our parameter fusion strategy effectively eliminates this overhead, maintaining the exact same computational efficiency as the original backbone. Implementation details are provided in Section~G.

    \textit{Evidence from metrics:}  
    Relative to the baseline \textit{DenoiseRep} module, multi-step fusion further enhances feature refinement 
    and separation quality. This results in performance improvements such as COCO AP increasing 
    from 44.3\% to 44.7\% (\textbf{+0.40\% over \textit{DenoiseRep}}) and ADE20K mIoU rising from 34.8\% 
    to 35.0\% (\textbf{+0.20\% over \textit{DenoiseRep}}).
\end{itemize}


% --------------------------------
% 2. 实验分析
% --------------------------------
\vspace{0.6em}
\textbf{2. Experimental Analyses}

\begin{itemize}

    \item \textbf{Ablation Studies:}  
    Each component—teacher–student distillation and multi-step fusion—contributes consistently 
    across ImageNet classification, person ReID, COCO detection, and ADE20K segmentation.  
    The detailed experimental results are reported in \textbf{Tables~I} and \textbf{V}.

    \item \textbf{Feature Statistics Visualization:}  
    t-SNE plots, intra/inter-class distance statistics, and overlap measurements collectively show that 
    \textit{DenoiseRep\textsuperscript{++}} forms tighter clusters and improves inter-class separability.  
    Quantitative metrics (lower intra variance, higher inter mean) validate the effectiveness of the denoising module.  
    Full results are presented in \textbf{Table~XI}, and the corresponding visualization is shown in Figures~4 and~5.

\end{itemize}


% --------------------------------
% 3. 性能提升
% --------------------------------
\vspace{0.6em}
\textbf{3. Performance Improvements}

The improved system delivers consistent gains across multiple benchmarks:

\begin{center} \begin{tabular}{lcccccc} \toprule \textbf{Task} & \textbf{Dataset} & \textbf{Backbone} & \textbf{Baseline} & \textbf{+\textit{DenoiseRep}} & \textbf{+\textit{DenoiseRep\textsuperscript{++}}} \\ \midrule Classification & ImageNet-1k & Swin-T & 81.82\% & 82.13\% & \textbf{82.30\%} \\ Person Re-ID & MSMT17 & ViT-S & 66.30\% & 67.33\% & \textbf{67.73\%} \\ Detection & COCO & Swin-T & 42.8\% & 44.3\% & \textbf{44.7\%} \\ Segmentation & ADE20K & ResNet-50 & 34.1\% & 34.8\% & \textbf{35.0\%} \\ \bottomrule \end{tabular} \end{center}

\textbf{These gains are achieved with no additional inference cost},  
as all auxiliary computations are fused into a single lightweight module.\\[1.0em]


\hspace*{2em}
Thank you for considering our manuscript for publication in \textit{TPAMI}. 
We look forward to your feedback and hope that our work contributes to advancing 
the field of pattern analysis and machine intelligence.


% \let\newpage\relax
\bibliographystyle{ieeetr}
\bibliography{reference}
\end{document}
